\honest{The Generalized Anti-Zombie Principle}{Eliezer Yudkowsky}{2008}

I open with Descartes: each problem solved became a rule for solving others. \nb{Descartes wrote of truths found, not problems solved, in a passage on geometry and algebra. The sense is the same.}

Zombies are beings atom-for-atom identical to us, but not conscious. The core argument against them is simple. When you attend to your awareness, your inner voice says ``I am aware of being aware,'' and then you may say it aloud and type it into a blog post. So consciousness is not epiphenomenal. \nb{Chalmers held that the zombie argument does not require epiphenomenalism; see the notes on ``Zombie Responses.''}

I admit the question is ``philosophically controversial,'' but ``it seems like a considerable majority of Overcoming Bias commenters do accept this.'' Some conclusions can be reached only after accepting that you cannot subtract consciousness and leave the universe the same. So I ask you to accept it and move on.

How far does the argument generalize? Suppose someone has a switch that does not affect your brain, and says that flipping it will end your consciousness. You would go on talking about consciousness ``for exactly the same reasons,'' so if you are conscious now you would still be conscious.

This does not equate consciousness with talk about consciousness. Consciousness is, ``among other things,'' the cause of that talk in humans. A definition need not give necessary and sufficient conditions; sometimes it is a treasure map to the referent. If what in fact makes me talk about an unspeakable awareness is not consciousness, discourse becomes futile. That is no knockdown argument, since difficulties of discourse settle no empirical question, but whoever defies the principle has trouble with the meaning of their words, not only their plausibility. I could define the word as ``whatever actually makes humans talk about `consciousness','' and that would guarantee the word names something real. Even if consciousness is a confusion, the word would name the cognitive architecture that produces the confusion. But a definition is only a promise to use a word consistently, and settles no empirical question.

So the principle cannot say only that no change undetectable in principle can remove consciousness: the switch's flip is detectable, and still highly unlikely to remove it. Nor can it say that no change that leaves you undetectably affected can, because the switch does affect you. Its gravity pulls on your atoms, about $6 \times 10^{-16}$ m/s$^2$ from ten metres. \nb{The figure checks, and ``around half a neutron diameter per second per second'' is close to the true 0.4.} Flipping it light-years away would avoid that, but we should not have to change the thought experiment: whatever you are, you do not expect a disconnected switch across the room to touch it. That is a large step, though if you deny it you had better never go near a switch again. Physics is single-level, and the flip moves your particles by whole neutron diameters. What stays the same is only the higher-level description, of cells, proteins and spikes, since the map assigns many states one description. By the molecular level that force is no longer tracked.

But a small effect is not no effect. Enough tiny pulls would tear you apart, and by an amazing coincidence the pull could move one calcium ion slightly closer to a channel, making a neuron fire a bit sooner, a difference that amplifies chaotically into a different train of thought, an epileptic fit, and death. So I retreat: the pull is far below thermal noise, and thermal noise does not switch consciousness off. If you do not expect consciousness to flicker with thermal jiggling, you should not expect it to go out when someone sneezes a kilometre away.

My argument is now about expectations, not certainty, though the laws of probability are theorems, not suggestions. It lacks the clean form of ``You can't possibly eliminate consciousness while leaving all the atoms in exactly the same place.'' But it still carries: I don't know what consciousness is, but whatever causes my talk about it happens in my skull, in neurons or perhaps microtubules or neurotransmitters, and the switch affects those much less than thermal noise at 310 Kelvin, so I expect to talk about consciousness in almost exactly the same way afterward. It is weaker, I admit, but more general, and ``in terms of sheer common sense, correct.''

A reductionist and a substance dualist can both accept it. The reductionist thinks the important parts happen at a functional level far above atomic nuclei, so that someone who understood consciousness could describe it in terms of cognitive architecture. The substance dualist thinks it may involve quantum effects, but that if it flickered with every sneeze we would notice, like skipping a few seconds or waking from anaesthesia, and sometimes say ``I don't think therefore I'm not.'' Thermal noise does not disturb it, so the switch will not.

Hearing someone say ``Abracadabra'' does change your brain noticeably, and you may wonder why they said it, but you still go on talking about consciousness for almost the same reasons. Consciousness is not equated with that talk, but it is among the causes of it, so anything that turned consciousness off should stop the talk. The principle: if something cannot change your internal narrative, it should not make you cease to be conscious, and this holds while the narrative and its causes are ``pretty much the same.''

Why go to such lengths? Because of this debate. Albert proposes replacing your neurons with tiny robots with the same connections, input-output behaviour, internal state and learning rules. Bernice says that would kill you. Charles says there would be a conscious person, but not you. Sir Roger Penrose says the replacement is impossible without quantum gravity, and wanders away. Albert adds that the swap happens one neuron at a time: a robot surrounds a neuron, scans it, learns to copy it and takes over between one spike and the next, so that your behaviour changes by much less than thermal noise. Your inner narrative is unchanged, so does the principle not apply? \nb{The replacement case is older than the post; Chalmers lists versions by Pylyshyn (1980), Cuda (1985) and Searle (1992). Searle described Bernice's outcome, experience shrinking ``to nothing, while your externally observable behavior remains the same.'' Penrose's line matches Penrose's published view that consciousness depends on non-computable physics.}

Albert also argues that he need not even tell you: if your introspective evidence that you are the person of five minutes ago is unchanged, your conclusion is equally justified, and the principle says a change in your consciousness, let alone your identity, is something you ought to notice. \nb{This is the argument Chalmers made in 1995 for Albert's side on consciousness, with a switch between a neural and a silicon circuit: ``whenever experiences change significantly and we are paying attention, we can notice the change.'' The zombie argument's main advocate holds Albert's view on this question.} Bernice answers that the robots are a detectable Zombie Master, and that once she is replaced there is no one to notice. Charles says the robots are not faking anything, only doing what neurons do in silicon instead of carbon, so the new person is conscious. But the operation replaced the true cause of Charles's talk with a different cause, which happens to be conscious too, and introspection is imperfect anyway. When Albert accuses them of positing epiphenomenal facts, both say they can detect the difference experimentally. Albert answers that he can detect the switch flipping too, and that what they detect makes no noticeable difference to the true cause of their talk. Charles replies that two people who talk about personal identity for similar reasons are not thereby the same person.

I think the principle supports Albert, but the reasons must wait for future posts, because this one is already too long. \nb{As stated here, the principle covers changes that leave the causes of the narrative ``pretty much the same,'' and the robots replace them, which is Charles's objection. The reasons came in ``Identity Isn't In Specific Atoms'' and ``Timeless Identity,'' neither of which is in this book.} How far the argument generalizes matters: the makeup of future galactic civilizations may depend on the answer.
